\documentclass[12pt, a4paper]{article}

\usepackage[utf8]{inputenc}
\usepackage{tempora}
\usepackage{geometry}
\usepackage{hyperref}
\usepackage{fancyvrb}
\usepackage{showexpl}
\usepackage{indentfirst}
\usepackage{listings}
\usepackage{xcolor}
\usepackage{fancyhdr}
\usepackage[T2A]{fontenc}
\usepackage[english]{babel}
\linespread{1.25}
\usepackage[skip=10pt plus1pt, indent=40pt]{parskip}

\definecolor{codegreen}{rgb}{0,0.6,0}
\definecolor{codegray}{rgb}{0.5,0.5,0.5}
\definecolor{codepurple}{rgb}{0.58,0,0.82}
\definecolor{backcolour}{rgb}{0.96,0.96,0.96}
\lstset {
    language=C,
    backgroundcolor=\color{backcolour},   
    commentstyle=\color{codegreen},
    keywordstyle=\color{magenta},
    numberstyle=\tiny\color{codegray},
    stringstyle=\color{codepurple},
    basicstyle=\ttfamily\footnotesize,
    breakatwhitespace=false,         
    breaklines=true,          
    captionpos=b,
    keepspaces=true,
    showspaces=false,
    showstringspaces=false,
    showtabs=false,
    tabsize=4
}

\title{A custom Hardware Abstraction Layer for the STM8S103F3P6 microcontroller}
\author{Lev Alyukov \\ \texttt{lev.alyukov@gmail.com}}
\date{}

\begin{document}
\maketitle

\section{Introduction}

In this documentation I explain the princtiple of my HAL. The HAL was\
developed from scratch in C90 and automated using a Makefile, formatting code using clang-format.\
I use the datasheets by STMicroelectronics and an original HAL as a reference.\
First of all, it's a study project, so I could write inefficient\
code but I will be glad if you submit a pull request with better code. Thanks!

For begin to working, you need to have: The STM8S103F3P6 microcontroller, ST-Link v2 debugger/programmer, 4 socket-to-socket jumper wires, Small Device C Compilier (SDCC), STM8SFLASH Software, GNU Make, Cppcheck and Git version control. If you have it --- well, we can move on the next steps.

\newpage

\subsection{Build and flashing of the MCU}

For such tasks I using a classic GNU Makefile:

\begin{lstlisting}
make <target>
\end{lstlisting}

Where the target is:

\begin{center}
  \begin{tabular}{|l|l|}
    \hline
    Target & Description\\
    \hline
    \texttt{all} & Syntax check of files, linking and compile, flashing received the HEX file.\\
    \hline
    \texttt{build} & Only linking and compile files only.\\
    \hline
    \texttt{flash} & Flash HEX file only.\\
    \hline
    \texttt{syntax} & Syntax check of files only.\\
    \hline
    \texttt{format} & Files formatting.\\
    \hline
    \texttt{clear} & Cleaning all an object files.\\
    \hline
    \texttt{misra} & Run the MISRA C:2012 standard.\\
    \hline
    \texttt{pdf} & Compilie a PDF file from TeX.\\
    \hline
  \end{tabular}
\end{center}

\end{document}
 
